\documentclass[11pt,a4paper]{article}

% ---------------------------------------------------------------- langue & fontes
\usepackage{fontspec}
\usepackage[french]{babel}
\frenchsetup{SmallCapsFigTabCaptions=false,ThinSpaceInFrenchNumbers=true}
\setmainfont{LibertinusSerif-Regular.otf}[
  BoldFont=LibertinusSerif-Bold.otf, ItalicFont=LibertinusSerif-Italic.otf,
  BoldItalicFont=LibertinusSerif-BoldItalic.otf, Numbers=OldStyle]
\setsansfont{Inter-Regular.otf}[
  BoldFont=Inter-Bold.otf, ItalicFont=Inter-Italic.otf,
  BoldItalicFont=Inter-BoldItalic.otf, Scale=0.94]
\setmonofont{FiraMono-Regular.otf}[BoldFont=FiraMono-Bold.otf, Scale=0.84]
\newfontfamily\tabfont{Inter-Regular.otf}[BoldFont=Inter-Bold.otf,
  ItalicFont=Inter-Italic.otf, BoldItalicFont=Inter-BoldItalic.otf,
  Scale=0.94, Numbers=Lining]

% ---------------------------------------------------------------- mise en page
\usepackage[a4paper,top=2.5cm,bottom=2.3cm,left=2.1cm,right=2.1cm,
            headheight=17pt,headsep=13pt,footskip=20pt]{geometry}
\usepackage{microtype}
\usepackage{xcolor}
\usepackage{booktabs,tabularx,longtable,array,colortbl,multirow}
\usepackage{titlesec,titletoc}
\usepackage{fancyhdr}
\usepackage{enumitem}
\usepackage[most]{tcolorbox}
\usepackage{fancyvrb}
\usepackage{pifont}
\usepackage{fontawesome5}
\usepackage{needspace}
\usepackage{ragged2e}
\usepackage[hidelinks]{hyperref}

% ---------------------------------------------------------------- palette
\definecolor{encre}   {HTML}{16181D}
\definecolor{accent}  {HTML}{1F4E5F}
\definecolor{accent2} {HTML}{8C4A2F}
\definecolor{gris}    {HTML}{6B7280}
\definecolor{filet}   {HTML}{D5DADE}
\definecolor{fondclair}{HTML}{F3F6F7}
\definecolor{fondcode}{HTML}{F7F7F4}
\definecolor{vert}    {HTML}{17693F}
\definecolor{rouge}   {HTML}{9B2C2C}
\definecolor{ambre}   {HTML}{A86A12}
\color{encre}

\hypersetup{colorlinks=true,linkcolor=accent,urlcolor=accent,
  pdftitle={Reunion de cadrage - Projet NLP 2},
  pdfsubject={Jalon J2 - choix des cas d'usage, cadrage, outils novateurs},
  pdfauthor={Equipe projet NLP 2 - M2 Data et IA, FGES}}

% ---------------------------------------------------------------- symboles
\newcommand{\okmark}   {\textcolor{vert}{\ding{51}}}
\newcommand{\komark}   {\textcolor{rouge}{\ding{55}}}
\newcommand{\warnmark} {\textcolor{ambre}{\faExclamationTriangle}}
\newcommand{\checkbox} {\textcolor{gris}{\ding{113}}}
\newcommand{\cutmark}  {\textcolor{accent2}{\ding{34}}}
\newcommand{\starmark} {\textcolor{ambre}{\ding{72}}}
\newcommand{\wipmark}  {\textcolor{ambre}{\faTools}}
\newcommand{\code}[1]  {{\ttfamily #1}}
\newcommand{\rulebreak}{}

% ---------------------------------------------------------------- titres
\titleformat{\section}[hang]
  {\sffamily\bfseries\LARGE\color{accent}}{\thesection}{13pt}{}
  [\vspace{2pt}{\color{filet}\titlerule[0.9pt]}]
\titlespacing*{\section}{0pt}{24pt plus 4pt minus 2pt}{11pt}
\titleformat{\subsection}
  {\sffamily\bfseries\large\color{encre}}{\thesubsection}{9pt}{}
\titlespacing*{\subsection}{0pt}{16pt plus 3pt minus 1pt}{6pt}
\titleformat{\subsubsection}
  {\sffamily\bfseries\normalsize\color{accent2}}{}{0pt}{}
\titlespacing*{\subsubsection}{0pt}{12pt plus 2pt minus 1pt}{4pt}
\setcounter{secnumdepth}{2}
\setcounter{tocdepth}{2}

% ---------------------------------------------------------------- en-têtes
\pagestyle{fancy}\fancyhf{}
% en-tête courant, redéfini par build.py selon le document
\newcommand{\entetecourant}{Projet NLP 2}
\fancyhead[L]{\sffamily\scriptsize\color{gris}\entetecourant}
\fancyhead[R]{\sffamily\scriptsize\color{gris}\thepage}
\renewcommand{\headrulewidth}{0pt}
\renewcommand{\headrule}{\vspace{-4pt}{\color{filet}\rule{\headwidth}{0.4pt}}}
\fancypagestyle{plain}{\fancyhf{}\renewcommand{\headrule}{}%
  \fancyfoot[C]{\sffamily\scriptsize\color{gris}\thepage}}

% ---------------------------------------------------------------- paragraphes & listes
\setlength{\parindent}{0pt}
\setlength{\parskip}{5.5pt plus 1.5pt minus 0.5pt}
\linespread{1.05}
\raggedbottom
\widowpenalty=10000 \clubpenalty=10000
\setlist[itemize]{leftmargin=1.15em,itemsep=2pt,topsep=3pt,parsep=0pt,label=\textcolor{accent}{\textbullet}}
\setlist[enumerate]{leftmargin=1.5em,itemsep=2.5pt,topsep=3pt,parsep=0pt,
  label=\textcolor{accent}{\sffamily\bfseries\arabic*.}}

% ---------------------------------------------------------------- tableaux
% \hspace{0pt} : glue initiale, sans quoi TeX refuse de couper le premier
% mot de la cellule — cause de débordements en colonne étroite.
\newcolumntype{L}[1]{>{\raggedright\arraybackslash
  \rightskip=0pt plus 1fil\relax\hspace{0pt}}p{#1}}
% \hspace{0pt} : sans cette glue, le whatsit de \color empêche TeX
% de couper le premier mot de la cellule (débordements en en-tête).
\newcommand{\thead}[1]{{\tabfont\bfseries\color{accent}\hspace{0pt}#1}}
\newenvironment{tabletype}
  {\par\addvspace{7pt}\begingroup
   \tabfont\setlength{\tabcolsep}{4pt}\renewcommand{\arraystretch}{1.26}
   \setlength{\LTleft}{0pt}\setlength{\LTright}{0pt}
   \setlength{\extrarowheight}{1pt}
   % élasticité infinie à droite : TeX ne coupe plus un mot pour « remplir »,
   % mais coupe encore lorsque c'est le seul moyen d'éviter un débordement.
   % césure interdite : les largeurs minimales sont mesurées, chaque mot tient.
   % Un débordement, s'il en restait un, serait signalé par le journal.
   \hyphenpenalty=9000\exhyphenpenalty=9000
   \emergencystretch=0pt\hbadness=10000\relax}
  {\endgroup\par\addvspace{9pt}}
\setlength{\aboverulesep}{0.5pt}\setlength{\belowrulesep}{0.6pt}
\arrayrulecolor{filet}
\setlength{\heavyrulewidth}{0.9pt}\setlength{\lightrulewidth}{0.45pt}

% ---------------------------------------------------------------- encadrés
\newtcolorbox{callout}{enhanced,breakable,
  colback=fondclair,colframe=accent,boxrule=0pt,leftrule=2.6pt,
  sharp corners,left=11pt,right=11pt,top=8pt,bottom=8pt,
  before skip=10pt,after skip=11pt,parbox=false}
\newtcolorbox{codebox}{enhanced,breakable,
  colback=fondcode,colframe=filet,boxrule=0.5pt,arc=2pt,
  left=9pt,right=9pt,top=7pt,bottom=7pt,
  before skip=10pt,after skip=11pt}

% ---------------------------------------------------------------- sommaire
\renewcommand{\contentsname}{Sommaire}
\titlecontents{section}[0pt]
  {\addvspace{4pt}\sffamily\bfseries\color{encre}}
  {\contentslabel[\thecontentslabel]{1.6em}}{}
  {\titlerule*[0.7pc]{\color{filet}.}\contentspage}
\titlecontents{subsection}[1.9em]
  {\sffamily\small\color{gris}}
  {\contentslabel[\thecontentslabel]{2.3em}}{}
  {\titlerule*[0.7pc]{\color{filet}.}\contentspage}

% ---------------------------------------------------------------- divers
\newlength{\metalab}
\newlength{\tw}\setlength{\tw}{\textwidth}
\newcommand{\settw}[1]{\setlength{\tw}{\dimexpr\textwidth-#1pt\relax}}
% tolérance relâchée : sans elle, TeX préfère déborder de quelques points
% plutôt que de desserrer une ligne contenant un long jeton en chasse fixe.
\emergencystretch=1.6em

% Document autonome : il reprend le préambule de l'équipe (fontes, palette, encadrés) et ajoute
% ce que la chaîne Markdown ne sait pas produire : formules, schémas TikZ et graphiques pgfplots.
% Compilation depuis ce dossier : latexmk -xelatex 01_encodeur_e5_small.tex

\usepackage{amsmath}
\usepackage{unicode-math}
\setmathfont{LibertinusMath-Regular.otf}
\usepackage{tikz}
\usetikzlibrary{arrows.meta,positioning,calc,fit,backgrounds,babel,shapes.geometric}
\usepackage{pgfplots}
\pgfplotsset{compat=1.18}

\hypersetup{pdftitle={Comprendre l'encodeur multilingual-e5-small},
  pdfsubject={Support d'apprentissage pour l'indexation du corpus},
  pdfauthor={Equipe Data Tigers - NLP 2 - M2 Data et IA, FGES}}
\renewcommand{\entetecourant}{Comprendre l'encodeur multilingual-e5-small \textbullet{} Équipe Data Tigers}

\newcolumntype{R}[1]{>{\raggedleft\arraybackslash\hspace{0pt}}p{#1}}
% colonne extensible avec glue initiale : sans elle, le changement de couleur de \thead fait passer
% l'en-tête à la ligne suivante (même mécanisme que celui décrit dans le préambule)
\newcolumntype{Y}{>{\raggedright\arraybackslash\hspace{0pt}}X}
\newcommand{\tok}[1]{\tikz[baseline=(t.base)]\node[draw=filet,fill=fondcode,rounded corners=1.5pt,
  inner xsep=2.2pt,inner ysep=1.6pt,font=\ttfamily\small](t){\strut #1};}
\newcommand{\sep}{\hspace{1.6pt plus 1pt}\allowbreak}
\newcommand{\debutmot}{\textcolor{accent2}{\_}}
\newcommand{\mesure}{\textcolor{vert}{\faRuler}\,}
\newcommand{\source}{\textcolor{accent}{\faBook}\,}

\tikzset{
  boite/.style={draw=accent,fill=fondclair,rounded corners=2pt,align=center,
    font=\sffamily\small,inner sep=5pt,minimum height=9mm},
  boite2/.style={boite,draw=accent2,fill=white},
  fleche/.style={-{Stealth[length=2.4mm]},draw=encre,line width=0.7pt},
  etiquette/.style={font=\sffamily\scriptsize,text=gris,align=center},
}

\addto\captionsfrench{\renewcommand{\contentsname}{Sommaire}}
\begin{document}
\thispagestyle{empty}
\begingroup\sffamily
{\color{accent}\rule{\textwidth}{2.6pt}}\par\vspace{13pt}
{\fontsize{27}{31}\selectfont\bfseries\color{encre}\raggedright Comprendre l'encodeur\\multilingual-e5-small\par}\vspace{7pt}
{\fontsize{13.5}{17}\selectfont\color{gris} Équipe Data Tigers \textbullet{} NLP~2 \textbullet{} M2 Data \& IA \textbullet{} FGES}\par\vspace{5pt}
{\fontsize{11}{14}\selectfont\color{accent2}\bfseries Support d'apprentissage pour l'indexation du corpus (notebook 02)}\par\vspace{12pt}
{\color{filet}\rule{\textwidth}{0.6pt}}\par\vspace{14pt}
\endgroup

\begingroup\setlength{\parskip}{0pt}
\settowidth{\metalab}{\tabfont\bfseries Vérifications}
\begin{tabular}{@{}>{\tabfont\bfseries\color{accent}}l@{\hspace{16pt}}p{\dimexpr\textwidth-\metalab-18pt\relax}@{}}
Objet & Expliquer comment l'encodeur transforme un texte en vecteur, ce qu'il faut respecter pour l'utiliser, et comment lire ses résultats \\[3.5pt]
Public & Toute l'équipe, sans prérequis au-delà du module NLP~2 \\[3.5pt]
Version & 27 septembre 2026, révisée après la correction du découpage M2 et l'exécution du notebook 02 \\[3.5pt]
Vérifications & Chaque chiffre est soit cité d'une source primaire (\source), soit mesuré sur notre corpus (\mesure) \\[3.5pt]
\end{tabular}\par\vspace{16pt}\endgroup

\begin{callout}
Les exemples de ce document ont été calculés le 27 septembre 2026 sur les 4\,240 passages du découpage M2, avec
le modèle \code{intfloat/multilingual-e5-small} en précision fp32. Ce sont des \textbf{illustrations} : quatre
requêtes écrites pour l'occasion ne mesurent pas la qualité d'un système. L'évaluation se fera au notebook 03, sur
un jeu de questions rédigé par des membres de l'équipe qui ne construisent pas l'index.
\end{callout}

\vspace{10pt}{\color{filet}\rule{\textwidth}{0.4pt}}\par\vspace{16pt}
\begingroup\setlength{\parskip}{0pt}\tableofcontents\endgroup
\clearpage

% =====================================================================================
\section{Le problème que l'encodeur résout}

Le cas A3 demande de retrouver, pour une question en langage naturel, les passages du corpus qui y répondent, de
les classer, et d'afficher pour chacun un score, la référence de sa section et un lien vers la fiche. Aucun texte
n'est généré : le système retrouve des passages existants.

Deux familles de méthodes savent le faire.

\begin{itemize}
\item La \textbf{recherche lexicale} (TF-IDF, BM25) compare des mots. Un passage remonte s'il contient les mots
  de la question, pondérés par leur rareté. Elle est rapide, transparente, et échoue dès que la question n'emploie
  pas les mots du texte.
\item La \textbf{recherche dense} compare des sens. Un encodeur transforme la question et chaque passage en un
  vecteur de nombres ; deux textes de sens proche ont des vecteurs proches, même sans mot commun.
\end{itemize}

\subsubsection{Un exemple sur notre corpus}

Pour la question familière \emph{«~Peut-on être viré quand on est en arrêt de travail ?~»}, le mot «~viré~»
n'apparaît dans aucun passage \mesure. Un TF-IDF construit sur les passages M2 renvoie en tête une retranscription
sur le métier d'inspecteur, un passage sur la cause réelle et sérieuse, et un passage sur des prélèvements
atmosphériques (scores 0,277, 0,255 et 0,217). L'encodeur renvoie trois passages consacrés à l'arrêt de travail et
à ses effets sur le contrat (section \ref{sec:exemples}).

Une question ne prouve rien. Elle montre seulement ce que chaque famille sait faire. La recherche lexicale reste la
référence à battre : le rapport J2 a fixé comme critère que l'encodeur ne fasse pas moins bien que BM25 sur le même
jeu de questions.

% =====================================================================================
\section{Représenter un texte par un vecteur}

\subsection{Représentation creuse et représentation dense}

\begin{tabletype}
\begin{tabularx}{\textwidth}{@{}L{3.2cm}YY@{}}
\toprule
& \thead{TF-IDF (creuse)} & \thead{Encodeur e5 (dense)} \\
\midrule
Dimension & une par mot du vocabulaire : 15\,314 sur nos passages \mesure & 384, fixée par le modèle \\
Valeurs & presque toutes nulles : 75 valeurs non nulles en médiane par passage, soit 0,49\,\% \mesure & toutes non nulles \\
Sens d'une dimension & le poids d'un mot précis & aucun, pris isolément \\
Mots absents & invisibles : un mot inconnu ne compte pas & découpés en morceaux connus, et interprétés dans leur contexte \\
Apprentissage & aucun, sinon le calcul des fréquences & appris sur environ un milliard de paires de textes \\
\bottomrule
\end{tabularx}
\end{tabletype}

\subsection{La similarité cosinus}

La proximité de deux vecteurs $\mathbf{u}$ et $\mathbf{v}$ se mesure par le cosinus de l'angle qu'ils forment :
\[
  \cos(\mathbf{u},\mathbf{v}) \;=\; \frac{\mathbf{u}\cdot\mathbf{v}}{\lVert\mathbf{u}\rVert\,\lVert\mathbf{v}\rVert}
  \;=\; \frac{\sum_{i=1}^{384} u_i\,v_i}{\sqrt{\sum_i u_i^2}\,\sqrt{\sum_i v_i^2}} .
\]
Il vaut 1 pour deux vecteurs de même direction, 0 pour deux vecteurs orthogonaux. Si les vecteurs sont
\textbf{normalisés}, c'est-à-dire de longueur 1, le cosinus se réduit au produit scalaire $\mathbf{u}\cdot\mathbf{v}$.
C'est pourquoi on normalise tous les vecteurs : chercher les passages les plus proches revient alors à multiplier
une matrice par un vecteur.

\begin{center}
\begin{tikzpicture}[scale=2.1]
  \draw[filet,line width=0.6pt] (0,0) circle (1);
  \draw[gris,->] (-1.25,0) -- (1.3,0); \draw[gris,->] (0,-1.15) -- (0,1.25);
  \coordinate (q) at (22:1); \coordinate (p1) at (48:1); \coordinate (p2) at (150:1);
  \draw[fleche,accent] (0,0) -- (q); \node[font=\sffamily\small,text=accent,anchor=west] at (22:1.04) {question $\mathbf{q}$};
  \draw[fleche,vert] (0,0) -- (p1); \node[font=\sffamily\small,text=vert,anchor=south west] at (48:1.03) {passage pertinent $\mathbf{p}_1$};
  \draw[fleche,rouge] (0,0) -- (p2); \node[font=\sffamily\small,text=rouge,anchor=south east] at (150:1.03) {passage éloigné $\mathbf{p}_2$};
  \draw[accent2] (22:0.55) arc (22:48:0.55); \node[font=\sffamily\scriptsize,text=accent2] at (35:0.68) {$\theta_1$};
  \draw[gris] (22:0.3) arc (22:150:0.3); \node[font=\sffamily\scriptsize,text=gris] at (100:0.42) {$\theta_2$};
  \node[etiquette,anchor=west] at (1.35,-0.6) {vecteurs de longueur 1\\sur la sphère unité\\$\cos\theta_1 > \cos\theta_2$};
\end{tikzpicture}\\[2pt]
{\sffamily\scriptsize\color{gris} Schéma en deux dimensions ; les vecteurs réels en ont 384.}
\end{center}

% =====================================================================================
\section{Du texte au vecteur : l'architecture du modèle}

\begin{center}
\begin{tikzpicture}[node distance=5.5mm]
  \node[boite2,text width=3.1cm] (txt) {\code{query: Le salarié démissionnaire...}};
  \node[boite,right=of txt,text width=2.5cm] (tk) {Tokeniseur\\Unigram\\250\,002 tokens};
  \node[boite,right=of tk,text width=2.6cm] (emb) {Table\\d'embeddings\\250\,037 $\times$ 384\\+ positions};
  \node[boite,below=9mm of emb,text width=2.6cm] (enc) {12 couches\\Transformer\\(attention, 12 têtes)};
  \node[boite,left=of enc,text width=2.5cm] (pool) {Moyenne des\\tokens réels\\(masque)};
  \node[boite,left=of pool,text width=3.1cm] (norm) {Normalisation L2\\vecteur de 384\\nombres, longueur 1};
  \draw[fleche] (txt) -- (tk); \draw[fleche] (tk) -- (emb); \draw[fleche] (emb) -- (enc);
  \draw[fleche] (enc) -- (pool); \draw[fleche] (pool) -- (norm);
  \node[etiquette,above=1mm of tk] {jusqu'à 512 tokens};
  \node[etiquette,below=1mm of enc] {un vecteur de 384 par token};
  \node[etiquette,below=1mm of norm] {un seul vecteur par texte};
\end{tikzpicture}
\end{center}

\subsection{La tokenisation}

Le modèle ne lit pas des mots mais des \textbf{tokens}, morceaux de texte tirés d'un vocabulaire de 250\,002 unités
construit par l'algorithme Unigram de SentencePiece \source. Les mots fréquents sont un seul token, les mots rares
sont découpés en morceaux connus. Voici la tokenisation réelle de deux entrées \mesure ; le trait souligné
\debutmot{} marque un début de mot.

\begin{center}\begin{minipage}{\textwidth}\centering
\setlength{\fboxsep}{0pt}
\tok{<s>}\sep\tok{\debutmot que}\sep\tok{ry}\sep\tok{:}\sep\tok{\debutmot Le}\sep\tok{\debutmot salarié}\sep\tok{\debutmot d}\sep\tok{émission}\sep\tok{naire}\sep\tok{\debutmot per}\sep\tok{ço}\sep\tok{it}\sep\tok{-}\sep\tok{il}\sep\tok{\debutmot l}\sep\tok{'}\sep\tok{al}\sep\tok{location}\sep\tok{\debutmot chômage}\sep\tok{\debutmot ?}\sep\tok{</s>}\\[5pt]
\tok{<s>}\sep\tok{\debutmot passage}\sep\tok{:}\sep\tok{\debutmot Article}\sep\tok{\debutmot L}\sep\tok{.}\sep\tok{\debutmot}\sep\tok{1234}\sep\tok{-9}\sep\tok{\debutmot du}\sep\tok{\debutmot Code}\sep\tok{\debutmot du}\sep\tok{\debutmot travail}\sep\tok{</s>}
\end{minipage}\end{center}

Trois observations utiles.
\begin{itemize}
\item «~démissionnaire~» devient \code{d}, \code{émission}, \code{naire} : le modèle reconstruit le sens d'un mot
  rare à partir de ses morceaux et du contexte. Seul un caractère absent du vocabulaire deviendrait le
  token inconnu \code{<unk>} : il n'y en a aucun dans nos 4\,240 passages \mesure.
\item Le préfixe \code{query:} coûte 3 tokens, le préfixe \code{passage:} en coûte 2 \mesure. Ces tokens comptent
  dans la limite de 512 et participent à la moyenne finale.
\item Chaque texte est encadré par les tokens spéciaux \code{<s>} et \code{</s>}, eux aussi comptés dans la limite.
\end{itemize}

\subsection{La table d'embeddings et les positions}

Chaque token est remplacé par une ligne d'une table de 250\,037 lignes et 384 colonnes, à laquelle on ajoute un
vecteur de position (512 positions au plus). La table compte 35 lignes de plus que le vocabulaire du tokeniseur :
ces lignes ne sont jamais utilisées \mesure.

\begin{tabletype}
\begin{tabularx}{\textwidth}{@{}YR{3.2cm}R{2.2cm}@{}}
\toprule
\thead{Composant} & \thead{Paramètres} & \thead{Part} \\
\midrule
Table d'embeddings des tokens (250\,037 $\times$ 384) & 96\,014\,208 & 81,6\,\% \\
12 couches Transformer & 21\,293\,568 & 18,1\,\% \\
\quad dont, par couche, l'attention & 592\,128 & \\
\quad dont, par couche, le réseau à propagation avant & 1\,182\,336 & \\
Positions (512 $\times$ 384) et types de segment & 197\,376 & 0,2\,\% \\
Couche de sortie «~pooler~», inutilisée ici & 147\,840 & 0,1\,\% \\
Normalisation des embeddings & 768 & \\
\midrule
\textbf{Total} & \textbf{117\,653\,760} & 100\,\% \\
\bottomrule
\end{tabularx}
\end{tabletype}
{\sffamily\scriptsize\color{gris} Paramètres comptés sur le modèle chargé \mesure ; architecture lue dans son fichier de configuration \source.}

Plus des quatre cinquièmes du modèle sont la table d'embeddings, qui couvre une centaine de langues. Nos passages
n'en utilisent que 10\,081 lignes sur 250\,002, soit 4,0\,\% \mesure. C'est ce qui rend possible la réduction du
vocabulaire envisagée au jalon J2 pour alléger le modèle servi.

\subsection{Les 12 couches Transformer}

Chaque couche transforme la suite de vecteurs en une nouvelle suite de vecteurs, en deux temps.

\begin{enumerate}
\item \textbf{L'attention} permet à chaque token de regarder tous les autres. Pour chaque token, le modèle calcule
  une requête $Q$, une clé $K$ et une valeur $V$ ; le nouveau vecteur est une moyenne des valeurs, pondérée par la
  ressemblance entre sa requête et les clés des autres :
  \[
    \operatorname{Attention}(Q,K,V) \;=\; \operatorname{softmax}\!\left(\frac{QK^{\top}}{\sqrt{d}}\right)V ,
    \qquad d = 32 .
  \]
  Le calcul est fait 12 fois en parallèle (12 «~têtes~» de 32 dimensions, $12 \times 32 = 384$), chaque tête pouvant
  suivre un type de relation différent.
\item \textbf{Le réseau à propagation avant} transforme chaque vecteur isolément : 384 vers 1\,536 dimensions,
  activation GELU, puis retour à 384.
\end{enumerate}
Des connexions résiduelles et des normalisations entourent ces deux étapes. À la sortie de la douzième couche,
chaque token a un vecteur de 384 nombres qui dépend de toute la phrase : le «~salarié~» d'une phrase sur le
licenciement n'a pas le même vecteur que celui d'une phrase sur les congés.

\subsection{De plusieurs vecteurs à un seul : moyenne et normalisation}

Pour obtenir un vecteur par texte, on fait la \textbf{moyenne des vecteurs des tokens réels}, en ignorant les
tokens de remplissage ajoutés quand on encode plusieurs textes de longueurs différentes en même temps. Avec
$\mathbf{h}_t$ le vecteur du token $t$ et $m_t$ valant 1 pour un token réel et 0 pour un remplissage :
\[
  \mathbf{e} \;=\; \frac{\sum_t m_t\,\mathbf{h}_t}{\sum_t m_t},
  \qquad
  \hat{\mathbf{e}} \;=\; \frac{\mathbf{e}}{\lVert\mathbf{e}\rVert} .
\]
Le masque rend le résultat indépendant du remplissage : un passage encodé seul ou au milieu d'un lot de passages
plus longs donne le même vecteur, à $8{,}8\times10^{-8}$ près \mesure. Après normalisation, les 4\,240 vecteurs
ont une longueur comprise entre $0{,}99999982$ et $1{,}0000001$ \mesure.

% =====================================================================================
\section{Comment le modèle a appris}

\subsection{Trois étapes}

\begin{enumerate}
\item \textbf{Initialisation} à partir de \code{microsoft/Multilingual-MiniLM-L12-H384}, un modèle multilingue
  compact obtenu par distillation d'un modèle plus grand \source.
\item \textbf{Pré-entraînement contrastif faiblement supervisé} sur environ un milliard de paires de textes
  multilingues, sans annotation humaine : titre et contenu de pages web, titre et contenu d'articles de presse,
  paires de traductions, \textbf{titre de section et passage de Wikipédia}, question et réponse de forums. Lots de
  32\,000 paires, 30\,000 pas d'entraînement \source.
\item \textbf{Ajustement fin supervisé} sur environ 1,6 million de paires annotées de recherche d'information :
  MS~MARCO, Natural Questions, TriviaQA, SQuAD, Mr.~TyDi, MIRACL, entre autres. Lots de 512, deux époques, avec des
  négatifs difficiles et une distillation depuis un modèle croisé plus précis \source.
\end{enumerate}

Les paires «~titre de section, passage~» du pré-entraînement ressemblent à nos entrées, qui commencent par
«~fiche > section~». C'est une raison de penser que l'en-tête aide, pas une preuve : le notebook 03 mesurera
son effet en comparant les passages avec et sans en-tête.

\subsection{La perte contrastive InfoNCE}

L'apprentissage rapproche le vecteur d'une question de celui de son passage et l'éloigne de ceux des autres passages
du même lot, qui servent de négatifs. Pour un lot de $n$ paires $(q_i, p_i)$ :
\[
  \mathcal{L} \;=\; -\frac{1}{n}\sum_{i=1}^{n}
  \log \frac{\exp\big(\cos(q_i,p_i)/\tau\big)}{\sum_{j=1}^{n}\exp\big(\cos(q_i,p_j)/\tau\big)},
  \qquad \tau = 0{,}01 .
\]
C'est une classification : parmi les $n$ passages du lot, retrouver le bon. La température $\tau$ règle la sévérité
de la comparaison \source.

\needspace{13\baselineskip}
\begin{center}
\begin{tikzpicture}[x=8.5mm,y=8.5mm]
  \foreach \i in {1,...,5} {
    \node[font=\sffamily\scriptsize,text=accent] at (0.3,-\i+0.5) {$q_{\i}$};
    \node[font=\sffamily\scriptsize,text=accent2] at (\i+0.5,0.35) {$p_{\i}$};
    \foreach \j in {1,...,5} {
      \ifnum\i=\j
        \fill[vert!55] (\j,-\i) rectangle (\j+1,-\i+1);
      \else
        \fill[rouge!12] (\j,-\i) rectangle (\j+1,-\i+1);
      \fi
      \draw[white,line width=1pt] (\j,-\i) rectangle (\j+1,-\i+1);
    }
  }
  \node[etiquette,anchor=west,text width=6.6cm] at (6.6,-2.5) {Chaque ligne est une question du lot, chaque colonne
    un passage. La perte pousse la case verte (la bonne paire) à dominer toute sa ligne : les autres passages du lot
    jouent le rôle de mauvaises réponses, gratuitement.};
\end{tikzpicture}
\end{center}

\subsection{Pourquoi deux préfixes}

Le même encodeur sert pour les questions et les passages. Les auteurs brisent la symétrie en ajoutant
\code{query:} devant les questions et \code{passage:} devant les passages, chacun suivi d'une espace. Ils observent que cette asymétrie
compte surtout quand le corpus contient des paraphrases de la question \source. La documentation du
modèle précise : \code{query:} et \code{passage:} pour une recherche, \code{query:} des deux côtés pour une
similarité entre textes de même nature, comme le découpage sémantique M3 \source.

\subsection{Une conséquence : des scores resserrés}

Avec une température aussi basse, seul l'ordre des scores compte pendant l'apprentissage, pas leur valeur. Les
cosinus se concentrent donc dans une bande étroite, un comportement que la documentation annonce entre 0,7 et
1,0 \source. Sur notre corpus, deux passages pris au hasard ont un cosinus médian de 0,834 ; 90\,\% des paires sont
entre 0,790 et 0,873 \mesure. Un score de 0,83 ne signifie donc \textbf{pas} «~assez proche~».

% =====================================================================================
\section{Le bi-encodeur et l'architecture de notre site}

\subsection{Bi-encodeur contre modèle croisé}

\begin{tabletype}
\begin{tabularx}{\textwidth}{@{}L{3.2cm}YY@{}}
\toprule
& \thead{Bi-encodeur (notre choix)} & \thead{Modèle croisé} \\
\midrule
Principe & question et passage encodés séparément, puis comparés par produit scalaire & question et passage lus ensemble, un score par paire \\
Passages & encodés une seule fois, hors ligne & relus à chaque question \\
Coût par question & un encodage et un produit matrice-vecteur & un passage complet dans le modèle pour chaque candidat \\
Précision & bonne & meilleure, mais inapplicable dans un navigateur à cette échelle \\
\bottomrule
\end{tabularx}
\end{tabletype}

Le bi-encodeur est ce qui rend le cas A3 possible sans serveur : tout le travail lourd se fait une fois, chez nous.

\subsection{Ce qui se passe hors ligne et en ligne}

\begin{center}
\begin{tikzpicture}[node distance=5mm and 6mm]
  \node[boite2,text width=2.6cm] (fiches) {Passages M2\\4\,240 textes};
  \node[boite,right=of fiches,text width=3.0cm] (encpy) {Encodeur fp32\\Python, PyTorch\\\code{passage:} + texte};
  \node[boite,right=of encpy,text width=3.0cm] (index) {Index statique\\matrice 4\,240 $\times$ 384\\+ métadonnées};
  \node[boite2,below=19mm of fiches,text width=2.6cm] (question) {Question\\du visiteur};
  \node[boite,right=of question,text width=3.0cm] (encjs) {Encodeur ONNX q8\\navigateur\\\code{query:} + question};
  \node[boite,right=of encjs,text width=3.0cm] (score) {Produit scalaire\\avec les 4\,240\\vecteurs};
  \node[boite2,right=of score,text width=2.4cm] (res) {Top $k$ : score,\\section, lien};
  \draw[fleche] (fiches) -- (encpy); \draw[fleche] (encpy) -- (index);
  \draw[fleche] (question) -- (encjs); \draw[fleche] (encjs) -- (score); \draw[fleche] (score) -- (res);
  \draw[fleche,dashed] (index) -- (score);
  \begin{scope}[on background layer]
    \node[fit=(fiches)(encpy)(index),fill=fondcode,rounded corners=3pt,inner sep=7pt,
      label={[etiquette,anchor=south west]north west:hors ligne, une fois (notebook 02)}] {};
    \node[fit=(question)(encjs)(score)(res),fill=fondcode,rounded corners=3pt,inner sep=7pt,
      label={[etiquette,anchor=north west]south west:en ligne, à chaque question, sur la machine du visiteur}] {};
  \end{scope}
\end{tikzpicture}
\end{center}

La recherche est \textbf{exhaustive} : la question est comparée aux 4\,240 vecteurs, soit environ 1,6 million de
multiplications, négligeables pour un navigateur. L'index pèse 6,21~Mio en fp32 et 1,55~Mio en entiers de 8 bits
\mesure. Aucune «~base vectorielle~» n'est nécessaire : ces outils et leurs index approchés servent à interroger des
corpus bien plus grands, sans parcourir tous les vecteurs ; la plupart fonctionnent comme un serveur, et le sujet exclut tout serveur d'inférence.

\subsection{Les fichiers du modèle et la quantification}

Le dépôt \code{Xenova/multilingual-e5-small} fournit le même modèle au format ONNX, en plusieurs précisions \source :

\begin{tabletype}
\begin{tabularx}{\textwidth}{@{}L{2.2cm}L{4.4cm}R{2.0cm}Y@{}}
\toprule
\thead{Type} & \thead{Fichier} & \thead{Taille} & \thead{Remarque} \\
\midrule
fp32 & \code{model.onnx} & 448,5~Mio & chargé par défaut hors WASM, y compris en WebGPU \\
fp16 & \code{model\_fp16.onnx} & 224,4~Mio & demi-précision \\
q8 & \code{model\_quantized.onnx} & 112,8~Mio & \textbf{chargé par défaut en WASM} \\
int8, uint8 & \code{model\_int8.onnx}, \code{model\_uint8.onnx} & 112,6~Mio & autres variantes 8 bits \\
q4 & \code{model\_q4.onnx} & 380,2~Mio & plus lourd que q8 \\
\bottomrule
\end{tabularx}
\end{tabletype}
{\sffamily\scriptsize\color{gris} Tailles lues sur le Hub \source ; types par défaut lus dans le code de Transformers.js 4.3.0, fichier \code{src/utils/dtypes.js} \source.}

Quantifier, c'est stocker chaque poids sur moins de bits : 8 au lieu de 32 divise la taille par quatre, au prix d'une
petite erreur sur chaque calcul. Le fichier q4 pèse plus que le q8 parce que, d'après les tailles, la quantification
4 bits ne touche que les couches de calcul et laisse la table d'embeddings en fp32 : $96{,}0$ millions de poids
de 4 octets font à eux seuls 366~Mio. Cette explication est déduite des tailles, pas vérifiée dans le graphe ONNX.

\begin{callout}
\textbf{Point de vigilance majeur, mesuré au notebook 02.} L'index est calculé en fp32 avec PyTorch, mais les
questions sont encodées dans le navigateur avec le modèle q8. Sur 822 requêtes de contrôle, le navigateur retrouve le
même premier résultat que le calcul exact pour 91,8\,\% d'entre elles, et ce premier résultat reste parmi les 5
premiers pour 821 requêtes sur 822 \mesure. Une simulation du q8 en Python annonçait 96,5\,\% : elle sous-estimait
l'écart, parce qu'ONNX Runtime Web n'exécute pas le fichier q8 exactement comme ONNX Runtime en Python. La mesure qui
fait foi se fait donc dans le navigateur. En fp32, le navigateur reproduit exactement le calcul de référence.
\end{callout}

% =====================================================================================
\section{Les contraintes à respecter}

\begin{tabletype}
\begin{tabularx}{\textwidth}{@{}L{3.6cm}YL{4.6cm}@{}}
\toprule
\thead{Contrainte} & \thead{Pourquoi} & \thead{Comment le vérifier} \\
\midrule
\code{passage:} devant chaque passage, \code{query:} devant chaque question & le modèle a appris ainsi ; sur nos quatre requêtes, mettre \code{passage:} devant la question fait tomber le premier résultat de référence aux rangs 1, 4, 5 et 9 \mesure & colonne \code{entree\_encodeur} de l'export, qui contient déjà le préfixe \\
512 tokens au plus, préfixe et tokens spéciaux compris & au-delà, la fin du texte est coupée sans avertissement & le notebook 01 garantit un maximum de 512 ; contrôler aussi la longueur des questions \\
Même tokeniseur des deux côtés & sinon les deux vecteurs ne parlent pas la même langue & tokenisations \code{intfloat} et \code{Xenova} identiques sur 500 blocs sur 500 \mesure \\
Moyenne des tokens et normalisation, des deux côtés & le modèle a été entraîné avec la moyenne & dans Transformers.js, options \code{pooling: "mean"} et \code{normalize: true} \\
Type de poids explicite dans le navigateur & le type par défaut dépend du périphérique : q8 en WASM, fp32 ailleurs & passer \code{dtype: "q8"} et vérifier les octets téléchargés \\
Alignement de l'index et des métadonnées & la ligne $i$ de la matrice doit correspondre au passage $i$ & empreinte commune dans le manifeste \\
Même texte que l'export & recomposer l'entrée risque de changer le préfixe ou l'en-tête & encoder \code{entree\_encodeur} tel quel \\
Versions figées & un autre modèle ou une autre révision donne d'autres vecteurs & révisions des dépôts et fichier \code{requirements-lock.txt} \\
\bottomrule
\end{tabularx}
\end{tabletype}

\begin{callout}
\textbf{Piège de la documentation.} L'exemple de la carte \code{Xenova/multilingual-e5-small} appelle l'extracteur
sans préfixe et sans option. Or, par défaut, Transformers.js ne fait \textbf{ni moyenne ni normalisation}
(\code{pooling: "none"}, \code{normalize: false}) \source : cet exemple renvoie un vecteur par token et non un
vecteur par texte. L'appel correct est :

\code{await extracteur("query: " + question, \{ pooling: "mean", normalize: true \})}
\end{callout}

% =====================================================================================
\needspace{10\baselineskip}
\section{Pourquoi nous ne retirons pas les mots vides}

Retirer les mots vides («~le~», «~de~», «~ne~», «~pas~», etc.) est une étape classique des méthodes par sac de mots
comme TF-IDF : ces mots, très fréquents, n'aident pas à distinguer les documents et bruitent le décompte. Nous ne
le faisons pas pour l'encodeur, pour quatre raisons.

\begin{enumerate}
\item \textbf{Le modèle a appris sur du texte naturel}, avec ses mots grammaticaux. Lui donner un texte amputé
  l'éloigne de ce qu'il a vu à l'entraînement, sur les passages comme sur les questions.
\item \textbf{La littérature va dans ce sens.} Dai et Callan (SIGIR~2019, tableau 3) retirent mots vides et ponctuation
  de requêtes de description, longues de 14 mots en moyenne : un modèle BERT perd en qualité (nDCG@20 de 0,529 à 0,503), tandis qu'un modèle par sac de mots
  gagne (de 0,404 à 0,427) \source. Leur modèle lit question et document ensemble, en anglais : c'est une indication
  pour notre encodeur, pas une preuve.
\item \textbf{Sur notre corpus, ces mots portent du droit} \mesure. La liste française du projet Snowball (154 mots)
  retire «~ne~», «~pas~», «~sans~», «~pour~», «~sur~» et «~par~» ; la liste stopwords-iso
  (691 mots) retire en plus «~plus~», «~moins~», «~avant~», «~après~», «~pendant~», «~sauf~», «~aucun~», «~peut~»
  et «~doit~». Or 28,3\,\% de nos passages contiennent une négation, 523 contiennent «~au moins~», 478 «~sans~».
  Avec l'une comme l'autre liste, «~Licenciement sans cause réelle et sérieuse~» et «~Licenciement pour cause
  réelle et sérieuse~» deviennent le même texte.
\item \textbf{Il n'y a rien à gagner.} Le retrait économiserait 33 à 39\,\% des tokens, mais aucun passage n'est
  tronqué : la limite de 512 est déjà respectée.
\end{enumerate}

Une liste construite sur la fréquence dans le corpus ne fait pas mieux : au-delà d'une fréquence documentaire de
0,3, elle retire «~employeur~», «~salarié~», «~salariés~», «~entreprise~», «~code~» et «~travail~» \mesure, qui
distinguent précisément «~l'employeur peut-il~» de «~le salarié peut-il~». Pour les références lexicales TF-IDF et
BM25, en revanche, la question se pose : le notebook 03 comparera BM25 sans liste, avec Snowball, avec Snowball
sans les négations, et avec une coupure par fréquence, avec et sans racinisation française.

% =====================================================================================
\section{Lire et interpréter les résultats}

\subsection{Un score n'est pas une probabilité}

Le graphique ci-dessous donne les scores des 4\,240 passages pour la question «~Quelles mentions sont interdites
sur le bulletin de paie ?~» \mesure. Le meilleur passage obtient 0,902, la médiane est à 0,801, et le plus mauvais
passage du corpus obtient encore 0,726.

\begin{center}\begin{minipage}{\textwidth}\centering
\begin{tikzpicture}
\begin{axis}[width=0.93\textwidth,height=5.6cm,ybar,bar width=0.0042,xmin=0.70,xmax=0.92,ymin=0,ymax=620,
  axis x line*=bottom,axis y line*=left,axis line style={gris},tick style={gris},
  xticklabel style={font=\sffamily\scriptsize,text=gris,/pgf/number format/.cd,use comma,fixed,fixed zerofill,precision=2},
  yticklabel style={font=\sffamily\scriptsize,text=gris},
  xlabel={\sffamily\scriptsize cosinus entre la question et le passage},ylabel={\sffamily\scriptsize passages},
  xlabel style={text=gris},ylabel style={text=gris},
  xtick={0.70,0.74,0.78,0.82,0.86,0.90},ytick={0,200,400,600},
  ymajorgrids,grid style={filet,line width=0.4pt},enlargelimits=false,clip=false]
\addplot[fill=accent!80,draw=none] coordinates {(0.7025,0) (0.7075,0) (0.7125,0) (0.7175,0) (0.7225,0) (0.7275,1) (0.7325,2) (0.7375,4) (0.7425,8) (0.7475,17) (0.7525,21) (0.7575,34) (0.7625,47) (0.7675,71) (0.7725,120) (0.7775,162) (0.7825,229) (0.7875,357) (0.7925,395) (0.7975,482) (0.8025,535) (0.8075,521) (0.8125,439) (0.8175,329) (0.8225,203) (0.8275,113) (0.8325,78) (0.8375,25) (0.8425,19) (0.8475,6) (0.8525,3) (0.8575,6) (0.8625,0) (0.8675,1) (0.8725,3) (0.8775,2) (0.8825,5) (0.8875,0) (0.8925,1) (0.8975,0) (0.9025,1) (0.9075,0) (0.9125,0) (0.9175,0)};
\draw[gris,dashed,line width=0.7pt] (axis cs:0.8015,0) -- (axis cs:0.8015,600);
\node[font=\sffamily\scriptsize,text=gris,anchor=south west] at (axis cs:0.8025,575) {médiane 0,801};
\draw[accent2,line width=1pt] (axis cs:0.9022,0) -- (axis cs:0.9022,250);
\node[font=\sffamily\scriptsize,text=accent2,anchor=south east,align=right] at (axis cs:0.9195,255)
  {meilleur passage 0,902\\7 passages au-dessus de 0,88};
\end{axis}
\end{tikzpicture}\\[-2pt]
{\sffamily\scriptsize\color{gris} Une barre par tranche de 0,005. Les passages utiles sont dans la queue de droite, presque invisible à cette échelle.}
\end{minipage}\end{center}

Trois règles de lecture en découlent.
\begin{itemize}
\item \textbf{Seul l'ordre compte.} Un score n'est comparable qu'aux autres scores de la même question.
\item \textbf{Il n'existe pas de seuil universel.} Pour masquer les résultats trop faibles ou décider une abstention
  (cas B1), un seuil doit être calibré sur le jeu de questions, et vérifié sur des questions hors corpus.
\item \textbf{Les écarts sont petits.} Entre le premier et le deuxième résultat, l'écart va de 0,0008 à 0,0095 sur
  nos quatre requêtes \mesure : le classement fin est fragile, d'où l'importance de mesurer l'écart fp32 et q8.
\end{itemize}

\subsection{Exemples réels, dont un échec}\label{sec:exemples}

\begin{tabletype}
\begin{tabularx}{\textwidth}{@{}R{1.0cm}R{1.35cm}L{5.1cm}Y@{}}
\toprule
\thead{Rang} & \thead{Score} & \thead{Fiche > section} & \thead{Début du passage} \\
\midrule
\multicolumn{4}{@{}l}{\textbf{«~Quelles mentions sont interdites sur le bulletin de paie ?~»}} \\
1 & 0,9022 & Le bulletin de paie > Et les mentions interdites ? & Aucune mention relative à l'exercice du droit de grève... \\
2 & 0,8927 & Le bulletin de paie > Chapô & Au moment du versement de son salaire, un bulletin de paie... \\
3 & 0,8839 & Le bulletin de paie > Quelles sont les mentions obligatoires ? & Les mentions suivantes doivent obligatoirement figurer... \\
\midrule
\multicolumn{4}{@{}l}{\textbf{«~Mon employeur peut-il me licencier pendant un arrêt maladie ?~»}} \\
1 & 0,8846 & Les absences liées à la maladie... > Peut-il y avoir licenciement pour maladie ? & L'état de santé du salarié ne peut pas constituer en soi un motif de licenciement... \\
2 & 0,8835 & Les absences liées à la maladie... > Chapô & En cas d'accident ou de maladie non professionnel, le salarié peut bénéficier d'un arrêt de travail... \\
\midrule
\multicolumn{4}{@{}l}{\textbf{«~Peut-on être viré quand on est en arrêt de travail ?~»}} \\
1 & 0,8764 & Les absences liées à la maladie... > Chapô & En cas d'accident ou de maladie non professionnel... \\
2 & 0,8755 & L'arrêt de travail pour accident du travail... > Que se passe-t-il pendant l'arrêt de travail ? & Le contrat de travail du salarié victime d'un accident du travail... \\
\midrule
\multicolumn{4}{@{}l}{\textbf{«~Combien de temps peut durer la période d'essai d'un CDI ?~»}} \\
1 & 0,8927 & Le contrat à durée déterminée (CDD) > Quelle est la durée de la période d'essai ? & Le contrat de travail à durée déterminée peut comporter une période d'essai... \\
5 & 0,8837 & La période d'essai > Quelle est la durée de la période d'essai ? & Au terme de l'article L. 1221-19 du Code du travail, le contrat de travail à durée indéterminée... \\
\bottomrule
\end{tabularx}
\end{tabletype}

La dernière requête est un \textbf{échec instructif}. Le passage qui répond, celui qui cite l'article L.~1221-19
propre au CDI, n'arrive qu'au rang 5, à 0,009 du premier, derrière un passage sur le CDD. Le passage CDD partage
avec la question l'expression «~période d'essai~» ; le passage CDI écrit «~contrat de travail à durée indéterminée~»
là où la question dit «~CDI~». Un essai va dans ce sens : en écrivant «~contrat à durée indéterminée~», ce passage
remonte au rang 2, derrière un autre passage de la même fiche \mesure. Une seule question ne prouve rien, mais le
jeu de questions devra contenir des sigles. Avec la métrique
Recall@5, cette question compte comme réussie ; avec Recall@1, comme ratée. Le choix de la métrique fait partie
de l'interprétation.

\subsection{Mesurer la qualité : les métriques}\label{sec:metriques}

Pour une question dont on connaît le ou les passages qui répondent :
\begin{itemize}
\item \textbf{Recall@}$k$ : 1 si un passage qui répond figure parmi les $k$ premiers, 0 sinon, en moyenne sur les
  questions. C'est la métrique la plus lisible pour un affichage de $k$ résultats.
\item \textbf{MRR} (rang réciproque moyen) : $\frac{1}{|Q|}\sum_{q} \frac{1}{\operatorname{rang}_q}$. Le premier bon
  résultat au rang 1 vaut 1, au rang 5 vaut 0,2.
\item \textbf{nDCG@}$k$ : récompense les bons résultats d'autant plus qu'ils sont hauts, quand plusieurs passages
  répondent à des degrés divers.
\end{itemize}
Deux méthodes se comparent sur les \textbf{mêmes questions}, par un test apparié de McNemar. D'après le calcul du
jalon J2, détecter un écart de 10 points avec un taux de désaccord de 0,20 demande environ 155 questions.

\subsection{Pièges d'interprétation}

\begin{itemize}
\item \textbf{Scores élevés partout} : deux passages sans rapport sont à 0,83 en médiane. Comparer des rangs, jamais
  des valeurs absolues.
\item \textbf{Doublons} : deux passages identiques ont le même vecteur et occupent deux places du classement. Le
  découpage M2 en compte 57 à texte identique, mais seulement 3 à entrée identique grâce à l'en-tête \mesure.
\item \textbf{Comparer deux modèles par leurs scores} n'a pas de sens : chaque modèle a sa propre échelle.
\item \textbf{Les projections en deux dimensions} (ACP, t-SNE, UMAP) aident à explorer mais déforment les distances :
  deux points voisins sur le dessin peuvent être éloignés dans les 384 dimensions, et inversement.
\item \textbf{Une réussite sur quelques questions} ne dit rien de la qualité générale : seules comptent les mesures
  sur le jeu de questions complet.
\end{itemize}

% =====================================================================================
\needspace{20\baselineskip}
\section{Pourquoi ce modèle}

Le choix a été fait au jalon J2 sous quatre contraintes : exécution dans le navigateur, coût nul, corpus en français,
poids téléchargé raisonnable. Sept candidats ont été mesurés (poids quantifiés en 8 bits, relevés au J2) :

\begin{tabletype}
\begin{tabularx}{\textwidth}{@{}L{9.0cm}R{2.0cm}Y@{}}
\toprule
\thead{Modèle} & \thead{Poids} & \thead{Verdict} \\
\midrule
\code{Xenova/multilingual-e5-small} & 112,8~Mio & retenu \\
\code{Xenova/paraphrase-multilingual-MiniLM-L12-v2} & 112,8~Mio & équivalent, sans avantage \\
\code{Xenova/distiluse-base-multilingual-cased-v2} & 129,0~Mio & plus lourd \\
\code{Xenova/multilingual-e5-base} & 265,7~Mio & 2,4 fois plus lourd \\
\code{onnx-community/embeddinggemma-300m-ONNX} & 294,6~Mio & hors budget \\
\code{Xenova/bge-m3} & 543,3~Mio & hors budget \\
\code{Xenova/all-MiniLM-L6-v2} & 21,9~Mio & anglais uniquement \\
\bottomrule
\end{tabularx}
\end{tabletype}

\subsubsection{Ce qui le justifie}
\begin{itemize}
\item \textbf{Entraîné pour la recherche}, avec des paires question-passage et des préfixes dédiés, contrairement
  aux modèles de paraphrase, entraînés à comparer des phrases de même nature.
\item \textbf{Multilingue, français compris} : 100 langues, et des résultats publiés en français sur MIRACL, un
  jeu de recherche sur Wikipédia : nDCG@10 de 47,6 pour la version small, 49,7 pour la base, 54,5 pour la large
  \source. Le rapport J2 affirmait qu'aucun résultat de recherche en français n'était publié : c'est inexact, et à
  corriger.
\item \textbf{Léger} : 384 dimensions pour un index de 6,21~Mio en fp32 ; 19,8~ms pour encoder une question dans
  le navigateur, mesurés au J2.
\item \textbf{Licence MIT}, compatible avec un site public.
\end{itemize}

\subsubsection{Ses limites}
\begin{itemize}
\item L'ajustement fin supervisé est surtout anglais, et aucune source d'entraînement n'est spécialisée en droit :
  le texte juridique n'y figure qu'au hasard des pages web. L'exemple CDD et CDI montre la conséquence possible
  sur un vocabulaire spécialisé.
\item Le fichier du modèle pèse 112,8~Mio (135,6~Mio téléchargés à la première visite, mesurés au J2), au-dessus
  de la limite de 100~Mio par fichier de GitHub : le modèle doit être servi depuis le Hub, ou réduit à son
  vocabulaire utile.
\item La limite de 512 tokens impose le découpage, et la moyenne des tokens peut diluer le sens d'un passage long.
\end{itemize}

Les deux voies d'amélioration déclarées au J2 répondent à ces limites : l'ajustement fin contrastif sur nos propres
paires question-passage, et la réduction du vocabulaire.

% =====================================================================================
\needspace{12\baselineskip}
\section{Glossaire}

\begin{tabletype}
\begin{tabularx}{\textwidth}{@{}L{3.6cm}Y@{}}
\toprule
\thead{Terme} & \thead{Définition} \\
\midrule
Token & morceau de texte du vocabulaire du modèle : un mot, une partie de mot ou un signe \\
Embedding, plongement & vecteur de nombres qui représente un token ou un texte \\
Encodeur & réseau qui transforme un texte en vecteurs \\
Bi-encodeur & encodeur qui traite question et passage séparément, pour les comparer ensuite \\
Moyenne masquée & moyenne des vecteurs des seuls tokens réels, sans les tokens de remplissage \\
Normalisation L2 & division d'un vecteur par sa longueur, pour qu'il soit de longueur 1 \\
Similarité cosinus & cosinus de l'angle entre deux vecteurs ; produit scalaire s'ils sont normalisés \\
Perte contrastive & fonction d'entraînement qui rapproche les bonnes paires et éloigne les mauvaises \\
Température & paramètre qui règle la sévérité de la perte contrastive \\
Quantification & stockage des poids sur moins de bits, pour réduire la taille du modèle \\
ONNX & format standard de modèle, exécutable dans le navigateur par ONNX Runtime Web \\
Recall@$k$, MRR, nDCG & métriques de qualité d'un classement, décrites en section \ref{sec:metriques} \\
\bottomrule
\end{tabularx}
\end{tabletype}

% =====================================================================================
\section{Sources}

\begingroup\raggedright\small
\begin{itemize}
\item Modèle \code{intfloat/multilingual-e5-small}, révision \code{614241f}, carte et fichiers de configuration :
  \url{https://huggingface.co/intfloat/multilingual-e5-small}
\item Conversion ONNX \code{Xenova/multilingual-e5-small}, révision \code{761b726} :
  \url{https://huggingface.co/Xenova/multilingual-e5-small}
\item L. Wang et al., \emph{Text Embeddings by Weakly-Supervised Contrastive Pre-training}, arXiv 2212.03533 :
  perte InfoNCE, température, préfixes.
\item L. Wang et al., \emph{Multilingual E5 Text Embeddings: A Technical Report}, arXiv 2402.05672 : données et
  hyperparamètres d'entraînement, résultats MIRACL par langue (tableau 6).
\item Z. Dai, J. Callan, \emph{Deeper Text Understanding for IR with Contextual Neural Language Modeling}, SIGIR~2019,
  arXiv 1905.09217, tableau 3.
\item Transformers.js 4.3.0, fichiers \code{src/utils/dtypes.js} et \code{src/pipelines/feature-extraction.js},
  paquet npm \code{@huggingface/transformers}.
\item Listes de mots vides : Snowball, \url{https://snowballstem.org/algorithms/french/stop.txt} ; stopwords-iso,
  \url{https://github.com/stopwords-iso/stopwords-fr}.
\item Mesures sur notre corpus : passages M2 du notebook \nolinkurl{scripts/decoupage/01_corpus_decoupage.ipynb}.
\end{itemize}
\endgroup

\end{document}
